% ============================================================================
% Fichier cible : ARIEL_LOGMINER_MEMOIRE_FINAL/chapters/chapitre2_etat_art.tex
%   (§1.24 « Positionnement scientifique d'Ariel Logminer », tableau 1.3
%   tab:positionnement-logminer, PDF p. 40)
% Passage à remplacer : les trois lignes du tableau qui commencent par
%   « Hétérogénéité & », « Mémoire & » et « Dashboard & » (chacune jusqu'à son
%   « \\ » final), et ajout d'une ligne « Priorisation » après « Mémoire ».
%   Les autres lignes sont inchangées.
% Auteur : Agent 2 (statut des revendications), chantier 2026-11-05.
% Motifs : absence de scénario hétérogène conjoint (L3, non-recevabilité) ;
%   priorisation jamais mesurée (L2 §2.1) ; H5 réfutée (gain de la mémoire) ;
%   utilité du tableau de bord non mesurée (L1 D.2). La colonne de gauche ne
%   revendique que ce que le corps établit.
% PRIORITÉ : l'Agent 3 réécrit l'état de l'art et le positionnement ; s'il
%   refond ce tableau, il reprend ces quatre lignes (contenu, non forme).
% ============================================================================

Hétérogénéité &
Lecture d'un ensemble documenté de formats et projection partielle dans un schéma commun ; chaque source est évaluée séparément. &
Compatibilité universelle avec tout format de journal ; analyse conjointe d'un même incident à partir de journaux système et de flux réseau. \\
Mémoire &
Persistance de l'historique et support d'une procédure de mise à jour contrôlée. &
Gain mesuré de la mémoire sur l'attribution des tâches ; apprentissage continu autonome en production. \\
Priorisation &
Mécanisme prévu par la conception. &
Amélioration mesurée de l'ordre de traitement des anomalies. \\
Dashboard &
Restitution fonctionnelle des informations et support de validation. &
Amélioration du travail de l'analyste démontrée par une étude utilisateur. \\
